\chapter{Introducción y Motivación}

\section{Contextualización del Proyecto}
La ciudad de Nueva York representa uno de los ecosistemas inmobiliarios y turísticos más densos, dinámicos y competitivos a nivel mundial. Con millones de visitantes anuales, plataformas de alojamiento compartido como Airbnb han revolucionado la manera en que se distribuye y consume el espacio habitable temporal. El conjunto de datos analizado en este trabajo, \textit{New York City Airbnb Open Data} de 2019 \cite{airbnb2019}, captura una radiografía transversal de este fenómeno, conteniendo información detallada de casi 50,000 listados activos.

\section{Identificación del Problema de Negocio}
Para una plataforma bilateral como Airbnb, mantener el equilibrio entre la oferta (anfitriones) y la demanda (huéspedes) es fundamental para la maximización de ingresos. Un problema persistente en este ecosistema es la **correcta valoración del precio nocturno** de las propiedades. Los anfitriones novatos a menudo carecen del conocimiento del mercado para fijar precios competitivos, lo que resulta en tasas de ocupación subóptimas (si el precio es excesivo) o pérdida de ganancias potenciales (si el precio es deficiente).

Este proyecto aborda la necesidad de crear un motor algorítmico capaz de estimar el valor intrínseco de un alojamiento basándose estrictamente en sus características observables (tipo de habitación, métricas de disponibilidad) y, críticamente, en su centralidad geográfica y ubicación.

\section{Objetivos del Proyecto}
El desarrollo técnico de este informe se estructuró para dar cumplimiento a tres objetivos analíticos principales:

\begin{enumerate}
    \item \textbf{Predicción de Precios en NYC (Supervisado):} Construir, optimizar y contrastar una batería extensa de modelos de regresión (desde modelos lineales penalizados hasta métodos de \textit{ensemblaje} avanzado como XGBoost \cite{chen2016xgboost} y Stacking) para predecir el precio del alojamiento de forma precisa y generalizable.
    \item \textbf{Inteligencia Espacial (Ingeniería de Variables):} Trascender la simple pertenencia a un vecindario mediante el cálculo de distancias reales a puntos neurálgicos de la ciudad, transformando metadatos de latitud y longitud en variables predictivas robustas.
    \item \textbf{Segmentación de Mercado (No Supervisado):} Implementar un proceso de agrupación algorítmica para descubrir submercados latentes dentro de NYC, perfilando la oferta más allá de divisiones administrativas, lo que permite detectar patrones de demanda y concentración de opciones \textit{premium} versus económicas.
\end{enumerate}

Para garantizar un rigor estadístico absoluto y reproducibilidad, el proyecto completo (desde el remuestreo inicial hasta el \textit{GridSearchCV}) ha sido anclado mediante un estricto control estocástico (\texttt{random\_state=42}).
